\section{Gauge Transformations}
\label{sec:gauge_transformations}

The scalar and vector potentials are not uniquely determined by the electric and magnetic fields. A gauge transformation changes the potentials while leaving all electromagnetic fields unchanged. This freedom is not an accidental feature of one convenient parametrization; it is a central structural property of electromagnetism.

\subsection{The gauge transformation}

Let \(\chi(\mathbf{x},t)\) be a sufficiently smooth scalar function. Define new potentials by

\[
\boxed{
\mathbf{A}'
=
\mathbf{A}
+
\nabla\chi
}
\]

and

\[
\boxed{
\Phi'
=
\Phi
-
\frac{\partial\chi}{\partial t}
}.
\]

The two potentials must transform together. Changing only \(\mathbf{A}\) or only \(\Phi\) generally changes the electric field.

\subsection{Magnetic-field invariance}

The transformed magnetic field is

\[
\begin{aligned}
\mathbf{B}'
&=
\nabla\times\mathbf{A}'\\
&=
\nabla\times
\left(
\mathbf{A}+
\nabla\chi
\right)\\
&=
\nabla\times\mathbf{A}
+
\nabla\times\nabla\chi.
\end{aligned}
\]

Since

\[
\nabla\times\nabla\chi=0,
\]

we obtain

\[
\boxed{
\mathbf{B}'=\mathbf{B}
}.
\]

\subsection{Electric-field invariance}

The transformed electric field is

\[
\begin{aligned}
\mathbf{E}'
&=
-\nabla\Phi'
-
\frac{\partial\mathbf{A}'}{\partial t}\\
&=
-\nabla
\left(
\Phi-
\frac{\partial\chi}{\partial t}
\right)
-
\frac{\partial}{\partial t}
\left(
\mathbf{A}+
\nabla\chi
\right).
\end{aligned}
\]

Expanding,

\[
\begin{aligned}
\mathbf{E}'
&=
-\nabla\Phi
+
\nabla
\frac{\partial\chi}{\partial t}
-
\frac{\partial\mathbf{A}}{\partial t}
-
\frac{\partial}{\partial t}
\nabla\chi.
\end{aligned}
\]

For smooth \(\chi\),

\[
\nabla
\frac{\partial\chi}{\partial t}
=
\frac{\partial}{\partial t}
\nabla\chi.
\]

The additional terms cancel, and therefore

\[
\boxed{
\mathbf{E}'=\mathbf{E}
}.
\]

\subsection{Gauge-equivalent potentials}

Two pairs

\[
(\Phi,\mathbf{A})
\]

and

\[
(\Phi',\mathbf{A}')
\]

are gauge equivalent if they are related by a function \(\chi\) through the transformation above.

The physical electromagnetic configuration is represented not by one pair of potentials but by an entire equivalence class:

\[
\boxed{
[(\Phi,\mathbf{A})]
=
\left\{
\left(
\Phi-
\frac{\partial\chi}{\partial t},
\mathbf{A}+
\nabla\chi
\right)
\right\}
}.
\]

Every member of the class produces the same \(\mathbf{E}\) and \(\mathbf{B}\).

\subsection{Example: two gauges for a uniform magnetic field}

For

\[
\mathbf{B}=B_0\mathbf{e}_z,
\]

the symmetric gauge is

\[
\mathbf{A}_{\mathrm{s}}
=
\frac{B_0}{2}
\left(
-y\mathbf{e}_x
+x\mathbf{e}_y
\right).
\]

The Landau-type gauge is

\[
\mathbf{A}_{\mathrm{L}}
=
B_0x\mathbf{e}_y.
\]

Their difference is

\[
\mathbf{A}_{\mathrm{L}}
-
\mathbf{A}_{\mathrm{s}}
=
\frac{B_0}{2}
\left(
y\mathbf{e}_x
+x\mathbf{e}_y
\right).
\]

Choose

\[
\chi(x,y)
=
\frac{1}{2}B_0xy.
\]

Then

\[
\nabla\chi
=
\frac{B_0}{2}
\left(
y\mathbf{e}_x
+x\mathbf{e}_y
\right),
\]

so

\[
\mathbf{A}_{\mathrm{L}}
=
\mathbf{A}_{\mathrm{s}}
+
\nabla\chi.
\]

The scalar potential is unchanged because \(\chi\) is time independent.

\subsection{Gauge fixing}

A gauge condition selects one representative, or a restricted set of representatives, from each gauge-equivalence class. Gauge fixing does not change the physical fields. It removes redundant descriptive freedom and can simplify the equations for the potentials.

Common gauge conditions include:

\begin{enumerate}
    \item the Coulomb gauge,
    \[
    \boxed{
    \nabla\cdot\mathbf{A}=0
    };
    \]

    \item the Lorenz gauge,
    \[
    \boxed{
    \nabla\cdot\mathbf{A}
    +
    \frac{1}{c^2}
    \frac{\partial\Phi}{\partial t}
    =0
    };
    \]

    \item the temporal gauge,
    \[
    \boxed{
    \Phi=0
    }.
    \]
\end{enumerate}

The Lorenz gauge is named after Ludvig Lorenz. It is Lorentz covariant, but its name is not ``Lorentz gauge.''

\subsection{Reaching the Coulomb gauge}

Suppose the original vector potential does not satisfy

\[
\nabla\cdot\mathbf{A}=0.
\]

After a gauge transformation,

\[
\nabla\cdot\mathbf{A}'
=
\nabla\cdot\mathbf{A}
+
\nabla^2\chi.
\]

To impose the Coulomb gauge, choose \(\chi\) to solve

\[
\boxed{
\nabla^2\chi
=
-
\nabla\cdot\mathbf{A}
}.
\]

The solvability and uniqueness of this Poisson equation depend on boundary conditions.

\subsection{Reaching the Lorenz gauge}

Under a gauge transformation,

\[
\begin{aligned}
&\nabla\cdot\mathbf{A}'
+
\frac{1}{c^2}
\frac{\partial\Phi'}{\partial t}
\\
&=
\nabla\cdot\mathbf{A}
+
\nabla^2\chi
+
\frac{1}{c^2}
\frac{\partial\Phi}{\partial t}
-
\frac{1}{c^2}
\frac{\partial^2\chi}{\partial t^2}.
\end{aligned}
\]

Using

\[
\Box
=
\frac{1}{c^2}
\frac{\partial^2}{\partial t^2}
-
\nabla^2,
\]

the new potentials satisfy the Lorenz condition when

\[
\boxed{
\Box\chi
=
\nabla\cdot\mathbf{A}
+
\frac{1}{c^2}
\frac{\partial\Phi}{\partial t}
}.
\]

This is a wave equation for the gauge function.

\subsection{Residual gauge freedom}

A gauge condition does not always determine the potentials uniquely.

In the Lorenz gauge, a further transformation preserves the condition when

\[
\boxed{
\Box\chi=0
}.
\]

This remaining freedom is called residual gauge freedom.

In the Coulomb gauge, residual transformations satisfy

\[
\boxed{
\nabla^2\chi=0
}
\]

together with the relevant boundary conditions.

\subsection{Temporal gauge}

To set

\[
\Phi'=0,
\]

choose \(\chi\) such that

\[
\frac{\partial\chi}{\partial t}
=
\Phi.
\]

One possible solution is

\[
\chi(\mathbf{x},t)
=
\int_{t_0}^{t}
\Phi(\mathbf{x},t')
\,\mathrm{d}t'
+
\chi_0(\mathbf{x}).
\]

The arbitrary spatial function \(\chi_0\) represents residual freedom.

\subsection{Gauge invariance and physical redundancy}

A gauge transformation does not map one physical electromagnetic state to a different physical state. It changes the mathematical representative of the same state.

This differs from a global physical symmetry. A spatial rotation can map a field configuration into a distinct rotated configuration. A gauge transformation leaves all local field strengths unchanged.

The distinction will become especially important when the electromagnetic field is coupled to charged matter.

\subsection{Boundary conditions and allowed gauge functions}

Not every function \(\chi\) is admissible in every problem. Gauge transformations must respect the boundary conditions, fall-off requirements, periodicity, and regularity assumptions imposed on the potentials.

For example, if the potentials are required to vanish sufficiently rapidly at spatial infinity, a gauge function that grows without bound may not preserve the allowed class of configurations.

Gauge transformations that vanish appropriately at the boundary are often called small gauge transformations. Transformations with nontrivial boundary behavior can carry additional global information.

\subsection{Four-dimensional preview}

With the convention

\[
A^{\mu}
=
\left(
\frac{\Phi}{c},
\mathbf{A}
\right),
\]

the gauge transformation becomes

\[
\boxed{
A'^{\mu}
=
A^{\mu}
-
\partial^{\mu}\chi
}.
\]

For the time component,

\[
\frac{\Phi'}{c}
=
\frac{\Phi}{c}
-
\frac{1}{c}
\frac{\partial\chi}{\partial t}.
\]

For the spatial components, since

\[
\partial^{i}=-\partial_i,
\]

we obtain

\[
\mathbf{A}'
=
\mathbf{A}+
\nabla\chi.
\]

The next section develops this covariant notation systematically.

\subsection{Common mistakes}

Typical errors include:

\begin{enumerate}
    \item transforming \(\mathbf{A}\) without transforming \(\Phi\);
    \item using the same sign in both three-dimensional transformation formulas;
    \item treating gauge-equivalent potentials as physically different electromagnetic fields;
    \item assuming that a gauge condition removes all residual freedom;
    \item confusing the Lorenz gauge with a Lorentz transformation;
    \item choosing a gauge function that violates the boundary conditions of the problem.
\end{enumerate}

\subsection{What has been established}

The transformations

\[
\mathbf{A}'
=
\mathbf{A}+
\nabla\chi,
\qquad
\Phi'
=
\Phi-
\frac{\partial\chi}{\partial t}
\]

leave \(\mathbf{E}\) and \(\mathbf{B}\) unchanged. Gauge conditions select convenient representatives but may leave residual freedom. The covariant expression

\[
A'^{\mu}
=
A^{\mu}-\partial^{\mu}\chi
\]

prepares the construction of the electromagnetic four-potential.
